\chapter{紧算子与谱定理}
\label{ch:b3:spectral}

对角化是有限维线性代数的
明珠：对称矩阵必有由特征向量组成的
标准正交基。在无穷维情形下，这一结论对有界
自伴算子一般不成立 --- 乘以 $x$ 的乘法算子
在 $L^2(\intcc01)$ 上 \emph{没有}任何特征值
（见\cref{exo:b3:spectral:6}） --- 但这一结论
以完美的形式对那些 \emph{几乎有限维
}的算子成立，即紧算子。紧自伴算子的谱定理
是应用泛函分析中最常用的定理：
它对角化积分方程，推动弗雷德霍姆择一性，
并在（周末问题）中求解弦振动问题，
产生傅里叶分析的正弦基，
完全来自纯算子理论 --- 并附带欧拉的
$\zeta(2) = \frac{\pi^2}6$ 由迹公式自然导出，
作为临别礼物。全文中，$H$ 是 $\C$ 上的希尔伯特空间（或
$\R$ 上的，结论作相应调整），算子均为有界的。

\section{紧算子}

\begin{definition}\label{def:b3:spectral:compact}
设 $T \in \mathcal L(E, F)$（$E, F$ 为巴拿赫空间），称其为
\emph{紧算子}\index{紧算子}，若单位球的像 $T(B)$
在 $F$ 中相对紧 --- 等价地，
每个有界列 $(x_n)$ 都有一个子列，使
$(Tx_{n_k})$ 收敛。有限秩算子都是紧算子
（有限维中的有界集是相对紧的）；无穷维空间的
恒等算子从来不是紧算子（里斯定理，本科二年级）。
\end{definition}

\begin{proposition}\label{prop:b3:spectral:ideal}
紧算子全体 $\mathcal K(E, F)$ 构成
$\mathcal L(E, F)$ 的闭子空间，并且是双边理想：若 $S$ 紧
$\Rightarrow$ $AS$ 与 $SB$ 均紧，其中 $A, B$ 为有界算子。
此外，在希尔伯特空间中，每个紧算子都是
有限秩算子的范数极限。
\end{proposition}

\begin{proof}
子空间：由序列刻画可直接得到。理想：
有界映射把收敛列变为收敛列，并
把有界集变为有界集。闭性：设 $T_n \to T$，且 $T_n$
均紧，$(x_k)$ 以 $1$ 为界；对角线抽取
使得 $(T_nx_{k_j})_j$ 对每个 $n$ 收敛；于是
$(Tx_{k_j})$ 为柯西列，因为
\[
\norm{Tx_{k_j} - Tx_{k_l}} \leq 2\vertiii{T - T_n} +
\norm{T_nx_{k_j} - T_nx_{k_l}} ,
\]
先选定 $n$，再选取下标。希尔伯特空间中的
逼近：设 $T$ 为紧算子，$K = \overline{T(B)}$ 为紧集；
给定 $\varepsilon$，用有限个球覆盖 $K$，
这些球为 $B(y_i, \varepsilon)$，记 $P$ 为
到 $V = \operatorname{Vect}(y_1, \dots, y_m)$ 的正交投影（该空间闭，
因其有限维）。于是 $PT$ 的秩有限，并且
对 $\norm x \leq 1$，取 $y_i$ 使 $\norm{Tx - y_i} <
\varepsilon$，则
\[
\norm{Tx - PTx} \leq \norm{Tx - y_i} + \norm{P(y_i - Tx)}
\leq 2\varepsilon
\]
（$y_i = Py_i$；$\vertiii P \leq 1$）：$\vertiii{T - PT} \leq
2\varepsilon$。
\end{proof}

\begin{example}\label{ex:b3:spectral:examples}
(a) $\ell^2$ 上的对角算子：$T(x_n) = (\lambda_nx_n)$
当且仅当 $\lambda_n \to 0$ 时为紧算子
（见\cref{exo:b3:spectral:2}）。
(b) $\mathcal C(\intcc01)$ 上的核算子：由阿斯科利定理，
它们为紧算子（见\cref{exo:b3:complete:7}）。
(c) \emph{希尔伯特–施密特算子}\index{希尔伯特–施密特
算子}：对 $k \in L^2(\intcc01^2)$，
\[
(T_kf)(x) = \int_0^1k(x, y)\,f(y)\,\dd y
\]
定义 $L^2(\intcc01)$ 上的一个紧算子，其范数满足
$\vertiii{T_k} \leq \norm k_{L^2}$
（见\cref{exo:b3:spectral:4}：截断
$k$ 的基展开，即可把 $T_k$ 表为有限秩算子的极限）。
\end{example}

\section{自伴算子}

\begin{definition}\label{def:b3:spectral:selfadjoint}
称 $T \in \mathcal L(H)$ 为\emph{自伴算子}\index{自伴
算子}，若 $T = T^*$（见\cref{exo:b3:hilbert:8}），即
$\langle Tx, y\rangle = \langle x, Ty\rangle$ 对所有 $x, y$ 成立。
于是 $\langle x, Tx\rangle \in \R$ 对每个 $x$ 成立（
即它等于其共轭）。
\end{definition}

\begin{proposition}\label{prop:b3:spectral:sanorm}
对自伴算子 $T$：\index{数值半径}
\[
\vertiii T = \sup_{\norm x \leq 1}\ \abs{\langle x,
Tx\rangle} .
\]
$T$ 的特征值都是实数，且不同特征值对应的
特征向量彼此正交。
\end{proposition}

\begin{proof}
记 $M$ 为上确界；$M \leq \vertiii T$ 由
柯西–施瓦茨不等式得出。反之，极化型恒等式
\[
\langle x{+}y, T(x{+}y)\rangle - \langle x{-}y,
T(x{-}y)\rangle = 4\operatorname{Re}\langle y, Tx\rangle
\]
（展开即可；交叉项 $\langle y, Tx\rangle + \langle x,
Ty\rangle = 2\operatorname{Re}\langle y, Tx\rangle$ 由自伴性），
结合平行四边形法则，得
\[
4\operatorname{Re}\langle y, Tx\rangle \leq
M\bigl(\norm{x{+}y}^2 + \norm{x{-}y}^2\bigr) = 2M\bigl(\norm
x^2 + \norm y^2\bigr).
\]
对 $\norm x = 1$ 且 $Tx \neq 0$，取 $y = Tx/\norm{Tx}$：
$4\norm{Tx} \leq 4M$。于是 $\vertiii T \leq M$。特征值：
若 $Tx = \lambda x$，$x \ne 0$，则 $\lambda\norm x^2 = \langle
x, Tx\rangle \in \R$。正交性：$\lambda\langle x,
y\rangle = \langle Tx, y\rangle = \langle x, Ty\rangle =
\mu\langle x, y\rangle$，其中 $\lambda \neq \mu$ 均为实数。
\end{proof}

\section{谱定理}

\begin{lemma}[存在极端特征值]
\label{lem:b3:spectral:existence}
设 $T \neq 0$ 为紧自伴算子，则 $\vertiii T$
或 $-\vertiii T$ 是 $T$ 的特征值。
\end{lemma}

\begin{proof}
由\cref{prop:b3:spectral:sanorm}，取单位向量 $x_n$
使 $\langle x_n, Tx_n\rangle \to \mu$，其中 $\abs\mu =
\vertiii T > 0$（取子列以固定符号）。则
\[
\norm{Tx_n - \mu x_n}^2
= \norm{Tx_n}^2 - 2\mu\langle x_n, Tx_n\rangle + \mu^2
\leq \vertiii T^2 - 2\mu\langle x_n, Tx_n\rangle + \mu^2
\longrightarrow 2\mu^2 - 2\mu\cdot\mu = 0 .
\]
由紧性，存在子列 $Tx_{n_k} \to y$；于是 $\mu
x_{n_k} = Tx_{n_k} - (Tx_{n_k} - \mu x_{n_k}) \to y$，故
$x_{n_k} \to x = y/\mu$，为单位向量，由连续性得
$Tx = \mu x$。
\end{proof}

\begin{theorem}[紧自伴算子的谱定理
]\label{thm:b3:spectral:spectral}
设 $T$ 为希尔伯特空间
$H$ 上的紧自伴算子。\index{谱定理（紧自伴算子）}
\begin{enumerate}
  \item $H$ 中存在标准正交系 $(e_n)_{n \in N}$
        （$N$ 为有限或可数），由 $T$ 的特征向量组成，
        其特征值 $(\lambda_n)$ 为非零实数，使得
        \[
        Tx = \sum_{n\in N}\lambda_n\,\langle e_n, x\rangle\,
        e_n \qquad (x \in H),
        \]
        且 $H = \ker T \,\oplus^\perp\,
        \overline{\operatorname{Vect}}(e_n : n \in N)$。
  \item 若 $N$ 为无穷集，则 $\lambda_n \to 0$；对每个
        $\delta > 0$，只有有限个 $n$ 满足
        $\abs{\lambda_n} \geq \delta$，且每个特征子空间
        $\ker(T - \lambda)$（$\lambda \neq 0$）是
        有限维的。
  \item 将 $(e_n)$ 用 $\ker T$ 的标准正交基补全，
        当 $H$ 可分时，所得为
        由特征向量组成的 $H$ 的标准正交基：$T$ 被对角化。
\end{enumerate}
\end{theorem}

\begin{proof}
先证 (2)。若有无穷多个标准正交特征向量 $x_k$
使 $\abs{\lambda_{(k)}} \geq \delta$，则 $\norm{Tx_k - Tx_l}^2 =
\lambda_{(k)}^2 + \lambda_{(l)}^2 \geq 2\delta^2$
（由正交性、毕达哥拉斯定理）：$(Tx_k)$ 无收敛子列，
与 $T$ 在有界列
$(x_k)$ 上的紧性矛盾。于是对每个 $\delta$，
特征值位于
$\intoo{-\delta}\delta$ 之外的总重数被一个有限数控制；可数性及 $\lambda_n \to 0$
随之成立。

(1) 设 $H_0$ 为由\emph{全部}
非零特征值的特征向量张成的闭线性包，按
(2) 及各有限维特征子空间内的格拉姆–施密特正交化，
以及不同特征子空间之间的正交性（由\cref{prop:b3:spectral:sanorm}），
整理为标准正交系
$(e_n)$，其特征值为 $\lambda_n \neq 0$。$T$ 将 $H_0$
映入 $H_0$，并且 $H_0^\perp$ 映入 $H_0^\perp$：对 $y
\perp H_0$，且 $e$ 为特征向量，$\langle e, Ty\rangle =
\langle Te, y\rangle = \lambda\langle e, y\rangle = 0$。
限制 $T' = T\restriction_{H_0^\perp}$ 为希尔伯特空间
$H_0^\perp$ 上的紧自伴算子；若 $T' \neq
0$，则由\cref{lem:b3:spectral:existence}得到
$T$ 的、特征值非零的特征向量，位于 $H_0^\perp$ 中 ---
这不可能，因为此类向量位于 $H_0$ 中。故 $T' = 0$：
$H_0^\perp \subseteq \ker T$。反之，$\ker T \perp$ 每个
$e_n$（$\langle e_n, z\rangle = \frac1{\lambda_n}\langle
Te_n, z\rangle = \frac1{\lambda_n}\langle e_n, Tz\rangle =
0$）：$\ker T \subseteq H_0^\perp$，故 $\ker T =
H_0^\perp$，从而得到正交分解。展开式为：
对 $x = z + \sum_nc_ne_n$（$z \in \ker T$，$c_n = \langle
e_n, x\rangle$；在 $H_0$ 上用\cref{thm:b3:hilbert:parseval}(1)），
由 $T$ 的连续性，得 $Tx = \sum_nc_n\lambda_ne_n$。

(3) 可分空间的闭子空间 $\ker T$ 仍可分，
即它有标准正交基
（\cref{prop:b3:hilbert:gramschmidt}）；它们的并为
$H$ 的标准正交基（由 (1) 中的分解）。
\end{proof}

\begin{theorem}[弗雷德霍姆择一定理]
\label{thm:b3:spectral:fredholm}
设 $T$ 为紧自伴算子，$\lambda \in \R\setminus
\{0\}$。\index{弗雷德霍姆择一定理}
\begin{enumerate}
  \item 若 $\lambda$ 不是特征值，则 $T - \lambda I$
        是双射且有有界逆：对每个 $f$，
        方程 $Tx - \lambda x = f$ 恰有一个
        解，且该解连续依赖于 $f$。
  \item 若 $\lambda$ 是特征值，则方程 $Tx - \lambda x = f$
        可解的充要条件是 $f \perp \ker(T - \lambda I)$，并且解
        在相差该（有限维）
        核中元素的意义下唯一。
\end{enumerate}
\end{theorem}

\begin{proof}
将 $x = z + \sum c_ne_n$ 和 $f = w + \sum d_ne_n$ 按
\cref{thm:b3:spectral:spectral} 分解（$z, w \in \ker T$）。方程
化为
\[
-\lambda z = w, \qquad (\lambda_n - \lambda)\,c_n = d_n\
(n \in N).
\]
(1) 若 $\lambda \notin \{\lambda_n\}\cup\{0\}$，由谱定理的
(2) 得 $\inf_n\abs{\lambda_n - \lambda} = \delta >
0$（特征值只在 $0 \neq \lambda$ 处聚集）。求解：
$z = -w/\lambda$，$c_n = d_n/(\lambda_n - \lambda)$，并且
$\sum\abs{c_n}^2 \leq \delta^{-2}\sum\abs{d_n}^2$，故存在唯一
解，满足 $\norm x \leq C\norm f$。
(2) 设 $\lambda = \lambda_{n}$（$n$ 属于有限集 $F$）：
方程 $(\lambda_n - \lambda)c_n = d_n$（$n \in F$）可解要求
$d_n = 0$，即 $f \perp e_n$（$n \in F$），亦即 $f
\perp \ker(T - \lambda I)$；此时 $c_n$（$n \in F$）
可任意取值。
\end{proof}

\begin{example}\label{ex:b3:spectral:minxy}
在 $L^2(\intcc01)$ 上取 $Tf(x) = \int_0^1\min(x, y)f(y)\dd y$：
一个核为实对称函数的希尔伯特–施密特算子：
它紧且自伴。求解 $Tf = \lambda f$：
关系式 $\bigl(Tf\bigr)(x) = \int_0^xyf(y)\dd y +
x\int_x^1f(y)\dd y$ 表明 $u = Tf$ 满足 $u'' = -f$（两次求导，
求导对连续的 $f$ 合法，且 $Tf$ 在
$f \in L^2$ 时连续，由控制收敛定理得），且 $u(0)
= 0$ 和 $u'(1) = 0$。于是特征函数满足 $\lambda u'' =
-u$，$u(0) = 0$，$u'(1) = 0$：
\[
u_n(x) = \sin\Bigl(\bigl(n + \tfrac12\bigr)\pi x\Bigr),
\qquad
\lambda_n = \frac{1}{\bigl(n + \frac12\bigr)^2\pi^2}
\quad (n \geq 0),
\]
而谱定理断言 --- 不借助傅里叶理论 ---
这些正弦函数在归一化后构成 $L^2(\intcc01)$ 的标准正交基
（$T$ 的核为 $0$：$Tf = 0$
由两次求导迫使 $f = 0$ 几乎处处成立）。
周末问题对振动弦沿用同一套思想，
并从迹中导出 $\zeta(2)$。
\end{example}

\begin{method}\label{met:b3:spectral:method}
给定积分方程或微分方程：(1) 将其改写为
$(I - \lambda K)u = f$ 或 $Ku = \lambda u$，其中 $K$ 为
积分算子；(2) 验证 $K$ 紧（希尔伯特–施密特核，
或阿斯科利定理），若可能，还验证其自伴（核为对称的
实函数）；(3) 用谱定理对角化，或
用弗雷德霍姆择一定理判断可解性；(4) 读出
解的存在性、唯一性、稳定性，以及
在特征基下的级数公式。微分算子是
无界的，但其 \emph{逆算子}（格林算子）
是紧的：总是先求逆。
\end{method}

\section{练习}

\begin{exercise}[$\star$]\label{exo:b3:spectral:1}
(a) 证明有限维值域的有界算子
是紧的。
(b) 证明赋范空间上的恒等算子紧，当且仅当
维数有限（里斯，本科二年级）。由此推出，
无限维空间上的紧算子永远不可逆，
即不存在有界逆。
\end{exercise}

\begin{exercise}[$\star$]\label{exo:b3:spectral:2}
设 $T(x_1, x_2, \dots) = (\lambda_1x_1, \lambda_2x_2, \dots)$
在 $\ell^2$ 上，且 $(\lambda_n)$ 有界。
(a) 证明 $\vertiii T = \sup\abs{\lambda_n}$。
(b) 证明 $T$ 紧，当且仅当 $\lambda_n \to 0$。
\emph{（对 $\Leftarrow$，截断；对 $\Rightarrow$，在
$(e_n)$ 上检验。）}
(c) 何时 $T$ 自伴？在该情形下用
直接观察验证谱定理。
\end{exercise}

\begin{exercise}[$\star\star$]\label{exo:b3:spectral:3}
给出理想性质的细节
（见\cref{prop:b3:spectral:ideal}）：若 $S$ 紧，$A, B$
有界，则 $ASB$ 紧。由此推出，若 $ST = TS = I$
对某有界 $S$，且 $\dim H = \infty$，则 $T$ 不是
紧的 --- 并与\cref{exo:b3:spectral:1}(b)协调。
\end{exercise}

\begin{exercise}[$\star\star$]\label{exo:b3:spectral:4}
（希尔伯特–施密特）设 $k \in L^2(\intcc01^2)$，$(e_n)$ 为
$L^2(\intcc01)$ 的希尔伯特基。
(a) 证明 $\vertiii{T_k} \leq \norm k_{L^2}$
\emph{（对 $y$ 用柯西–施瓦茨不等式，再用托内利定理）}。
(b) 将 $k(x,y) = \sum_{m,n}c_{mn}e_m(x)\overline{e_n(y)}$
在正方形上的 $L^2$ 中展开（证明这些乘积构成
该处的希尔伯特基），并证明截断该和式所得的
有限秩算子在算子范数下收敛于 $T_k$：
$T_k$ 是紧的。
\end{exercise}

\begin{exercise}[$\star\star$]\label{exo:b3:spectral:5}
设 $T$ 自伴，且 $\langle x, Tx\rangle \geq 0$ 对
所有 $x$ 成立（\emph{正}算子）。
(a) 证明特征值 $\geq 0$，且 $\vertiii T =
\sup_{\norm x\leq1}\langle x, Tx\rangle$。
(b) 证明广义柯西–施瓦茨不等式
$\abs{\langle x, Ty\rangle}^2 \leq \langle x, Tx\rangle\langle
y, Ty\rangle$
\end{exercise}

\begin{exercise}[$\star\star$]\label{exo:b3:spectral:6}
在 $L^2(\intcc01)$ 上，设 $(Mf)(x) = x\,f(x)$。
(a) 证明 $M$ 有界、自伴，且 $\vertiii M =
1$，但\emph{不存在}特征值。
(b) 证明 $M$ 不紧 \emph{（给出一个有界列，
其像不含收敛子列，例如
收缩到 $1$ 附近区间的归一化示性函数 --- 或
引用谱定理）}。
(c) \cref{lem:b3:spectral:existence}的证明在
$M$ 上何处失效？
\end{exercise}

\begin{exercise}[$\star\star$]\label{exo:b3:spectral:7}
（沃尔泰拉）在 $L^2(\intcc01)$ 上取 $Vf(x) = \int_0^xf(y)\dd
y$.
(a) 证明 $V$ 紧（为希尔伯特–施密特算子，核为 $\mathbf
1_{y < x}$），但不自伴；计算 $V^*$。
(b) 证明 $V$ 没有非零特征值。\emph{（由 $Vf =
\lambda f$：$f$ 有连续代表元，从而
是 $\mathcal C^1$ 函数，并满足 $\lambda f' = f$，$f(0) = 0$。）}
(c) 由此可得，仅有紧性不足以给出特征向量：
自伴性在\cref{thm:b3:spectral:spectral}中
是必不可少的。
\end{exercise}

\begin{exercise}[$\star\star\star$]\label{exo:b3:spectral:8}
（柯朗–费舍尔）设 $T$ 为紧、自伴、
\emph{正}的，特征值为 $\mu_1 \geq \mu_2 \geq \dots
> 0$（按重数计，特征向量为 $e_1, e_2,
\dots$）。证明：
\[
\mu_{k} = \max_{\substack{V \subseteq H \\ \dim V = k}}\
\min_{\substack{x \in V\\ \norm x = 1}}\ \langle x, Tx\rangle
= \min_{\substack{W \subseteq H\\ \operatorname{codim}W =
k-1}}\ \max_{\substack{x\in W\\ \norm x = 1}}\ \langle x,
Tx\rangle .
\]
\emph{（取 $V = \operatorname{Vect}(e_1,\dots,e_k)$；对上界，
任取 $V$，与
形如 $\operatorname{Vect}(e_k, e_{k+1}, \dots)$ 的空间求交：
由维数计数可知交集非零。）}由此推出
特征值关于 $T$ 单调（$T \leq S
\Rightarrow \mu_k(T) \leq \mu_k(S)$）。
\end{exercise}

\begin{exercise}[$\star\star$]\label{exo:b3:spectral:9}
利用 \cref{thm:b3:spectral:fredholm}，对 $Tf(x) =
\int_0^1\min(x,y)f(y)\dd y$（\cref{ex:b3:spectral:minxy}）：
问对于哪些 $\lambda \in \R$，积分方程
\[
f(x) - \lambda\int_0^1\min(x,y)\,f(y)\,\dd y = g(x)
\]
有唯一解 $f \in L^2$，且对每个 $g \in L^2$ 皆然？
例外值处的情形又如何？
\end{exercise}

\begin{exercise}[$\star\star\star$]\label{exo:b3:spectral:10}
设 $S$ 为 $\ell^2$ 上的移位算子
（见 \cref{exo:b3:banach:1}）。
(a) 证明 $S$ 没有特征值，而每个 $\lambda$
满足 $\abs\lambda < 1$，都是 $S^*$ 的特征值（显式
求出特征向量，它们是等比数列）。
(b) 验证 $S$ 与 $S^*$ 均不紧：
仿 \cref{exo:b3:spectral:2} 的做法，在 $(e_n)$ 上检验。
(c) 评注：对于非自伴、非紧的算子，
特征值的情形可以从空集一直到整个
圆盘 --- 真正保留下来的概念是\emph{谱}，
将在后续课程中研究。
\end{exercise}

\begin{exercise}[$\star\star$]\label{exo:b3:spectral:11}
（平方根）设 $T$ 为紧、自伴、正的算子，
（$\langle x, Tx\rangle \geq 0$），作用于希尔伯特空间 $H$ 上，
其谱分解为 $Tx = \sum_n\mu_n\langle e_n,
x\rangle e_n$（$\mu_n > 0$）。
(a) 定义 $Sx = \sum_n\sqrt{\mu_n}\,\langle e_n, x\rangle
e_n$；证明 $S$ 为紧、自伴、正算子，且 $S^2
= T$。
(b) 证明唯一性：任一紧、正、自伴的 $R$
满足 $R^2 = T$，则保持 $T$ 的各特征子空间
\emph{（$RT = R^3 = TR$：$R$ 与 $T$ 可交换，故
$R(\ker(T - \mu)) \subseteq \ker(T - \mu)$）}，而在
$\ker(T - \mu)$ 上，$R$ 为正算子，平方为
$\mu\,\mathrm{id}$（在有限维空间上）：
于该空间上对角化，并得出 $R =
\sqrt\mu\,\mathrm{id}$ 于每个特征子空间上，从而 $R = S$。
(c) 计算 $\sqrt G$，其中 $G$ 为
\cref{pb:b3:spectral:1} 中的弦算子，问哪个核在正弦基上的特征值为
$\frac1{n\pi}$？（将 $\sqrt G$ 表示为
核的 $L^2$ 极限；不要求给出闭形式。）
\end{exercise}

\begin{exercise}[$\star\star\star$]\label{exo:b3:spectral:12}
（奇异值分解）设 $T \in \mathcal L(H)$ 为
\emph{紧}算子，不必自伴。
(a) 证明 $T^*T$ 紧、自伴、正；设
$(e_n)$ 为标准正交的特征向量族，满足
$T^*Te_n = s_n^2e_n$，$s_n > 0$（称为\emph{奇异值}），
并以 $\ker(T^*T) = \ker T$ 的基补全（请证明此等式）。
(b) 设 $f_n = \frac{Te_n}{s_n}$，证明 $(f_n)$ 为
标准正交族，并建立如下的\emph{SVD}：
\[
Tx = \sum_n s_n\,\langle e_n, x\rangle\,f_n
\qquad (x \in H),
\]
级数在 $H$ 中收敛。
(c) 由此推出：$\vertiii T = \max_ns_n$；$T$ 为
有限秩算子的范数极限（重证
\cref{prop:b3:spectral:ideal} 在希尔伯特空间上的逆命题）；
对于 \cref{exo:b3:spectral:7} 中的沃尔泰拉算子 $V$，
它没有特征值，请说明为何奇异值分解仍然存在，
以及其各组成部分分别是什么
（将 $V^*V$ 识别为弦型核算子 ---
其特征值的显式计算
属于 \cref{exo:b3:spectral:9} 的范畴）。
\end{exercise}

\section{问题：弦振动与
\texorpdfstring{$\zeta(2)$}{zeta(2)}}

\begin{problem}[{周末问题 --- 格林算子、
正弦基与迹公式}]\label{pb:b3:spectral:1}
求解两端固定弦振动的特征值问题 ---
$-u'' = \nu u$，$u(0) = u(1) = 0$ ---
用算子理论，不做任何傅里叶计算即得正弦标准正交基，
并通过比较两种表达式来求 $\zeta(2)$，
这两种表达式都是格林算子的\emph{迹}。
在 $L^2(\intcc01)$ 上定义：
\[
(Gf)(x) = \int_0^1 g(x,y)\,f(y)\,\dd y,
\qquad
g(x, y) = \min(x,y)\,\bigl(1 - \max(x,y)\bigr).
\]

\medskip
\noindent\textbf{第一部分 --- 格林算子。}
\begin{enumerate}
  \item 证明 $g$ 连续、对称，且 $0 \leq g
        \leq \frac14$，$G$ 紧、
        自伴（\cref{ex:b3:spectral:examples}(c)）。
  \item 设 $f$ 连续，证明 $u = Gf$ 为 $\mathcal
        C^2$ 类函数，并满足
        \[
        -u'' = f, \qquad u(0) = u(1) = 0
        \]
        \emph{（将 $u(x) = (1-x)\int_0^xyf(y)\dd y +
        x\int_x^1(1-y)f(y)\dd y$，并两次求导）}。
        反之，若 $u \in \mathcal C^2$ 且 $u(0) = u(1)
        = 0$，则 $G(-u'') = u$：$G$ 是
        弦算子的逆。
  \item 证明 $\ker G = \{0\}$，\emph{（若 $Gf = 0$ 且 $f \in
        L^2$：取连续函数 $\varphi$ 作检验，并将 $G$
        借对称性与富比尼定理移到 $\varphi$ 上，再用
        \cref{cor:b3:lp:fundlemma} 中的基本引理 --- 或
        先作正则化）}，并证明 $G$ 为正算子：
        $\langle f, Gf\rangle \geq 0$。\emph{（对连续
        $f$：$\langle f, Gf\rangle = \int_0^1 (u')^2$，其中 $u
        = Gf$，分部积分，再由稠密性得结论。）}
\end{enumerate}

\medskip
\noindent\textbf{第 II 部分 --- 对角化：正弦
基。}
\begin{enumerate}[resume]
  \item 证明：$G$ 的特征函数，其
        特征值 $\lambda \ne 0$ 者，在相差常数倍的意义下，恰为
        $-\lambda u'' = u$，$u(0) = u(1) = 0$ 的解 \emph{（一个
        特征函数具有连续代表元 --- $Gf$
        连续（$f \in L^2$，为何？）--- 因此
        由第 2 题的自举法，它属于 $\mathcal C^2$）}。
  \item 解边值问题：
        $G$ 的特征值为 $\lambda_n = \frac1{n^2\pi^2}$（$n \geq 1$），
        归一化特征函数为 $e_n(x) =
        \sqrt2\,\sin(n\pi x)$；通过直接积分验证
        正交归一性，作为一项合理性检验。
  \item 由
        \cref{thm:b3:spectral:spectral} 及第 3 题可知，
        $\bigl(\sqrt2\sin(n\pi x)\bigr)_{n\geq1}$ 是一组
        标准正交\emph{基}（$L^2(\intcc01)$ 的）--- 无需
        斯通–魏尔斯特拉斯定理，也无需傅里叶级数。将
        $f(x) = x(1-x)$ 按此基展开，并写出
        它的帕塞瓦尔恒等式。
\end{enumerate}

\medskip
\noindent\textbf{第 III 部分 --- 迹公式与
$\zeta(2)$。}
\begin{enumerate}[resume]
  \item 证明下列两个恒等式
        \[
        \langle e_n, Ge_n\rangle = \lambda_n
        \quad\text{和}\quad
        \sum_{n\geq1}\lambda_n = \int_0^1 g(x,x)\,\dd x .
        \]
        对第二式（即\emph{迹公式}）：将
        $g(x, \cdot)$，在固定的 $x$ 下按基 $(e_n)$ 展开 ---
        证明其系数为 $\lambda_ne_n(x)$，于是
        $g(x, \cdot) = \sum_n\lambda_ne_n(x)\,e_n$ 在
        $L^2$ 中成立。此处正弦函数是显式的，可直接验证
        $\sum_n\lambda_ne_n(x)e_n(y)$ 收敛，
        \emph{一致}地在正方形上（与 $\sum
        \frac2{n^2\pi^2}$ 比较），故其和连续，且
        具有相同的 $L^2$ 展开（关于 $y$，对每个 $x$），因此等于
        $g(x,y)$ 处处成立。令 $y = x$，逐项
        积分。
  \item 计算 $\int_0^1g(x,x)\dd x = \int_0^1x(1-x)\dd x =
        \frac16$，由此得出
        \[
        \sum_{n\geq1}\frac{1}{n^2\pi^2} = \frac16,
        \qquad\text{即}\qquad
        \boxed{\ \zeta(2) = \frac{\pi^2}{6}\ } :
        \]
        由算子的迹得到的欧拉和。
  \item 用第三种方法重新推导 $\zeta(2)$：在
        正弦基中对常值函数 $\mathbf 1$ 应用帕塞瓦尔恒等式，
        计算 $\sum_{n \text{ 为奇数}}\frac1{n^2}$，并
        由此得出结论。随后比较两种机制：在何种意义上，
        第 7--8 题中的迹论证是“对整个核
        一次性应用帕塞瓦尔恒等式”？
\end{enumerate}

\medskip
\noindent\textbf{第 IV 部分 --- 弦的振动。}
\begin{enumerate}[resume]
  \item （综合运用分离变量法）设 $f \in
        L^2$，定义
        \[
        u(t, x) = \sum_{n\geq1}\;c_n\,
        \cos(n\pi t)\,\sqrt2\sin(n\pi x),
        \qquad c_n = \langle e_n, f\rangle .
        \]
        证明：该级数在 $L^2(\intcc01)$ 中对每个
        $t$ 收敛，$t\mapsto u(t, \cdot)$ 连续地映入
        $L^2$，并且当 $f$ 属于有限多个
        $e_n$ 张成的空间时，它满足波动方程 $\partial_t^2u =
        \partial_x^2u$，其中 $u(0) = f$，$\partial_tu(0) = 0$，
        端点固定。特征值 $n^2\pi^2$ 即平方频率，
        亦即弦的谐波 --- 请用一段话
        说明
        \cref{thm:b3:spectral:spectral} 的音乐解释。
\end{enumerate}

\medskip
\noindent\textbf{第 V 部分 --- 变分收益：幂法
、外尔稳定性与 $\pi$ 的严格上界。}令
$A$ 为紧自伴\emph{正}算子，其特征值为
$\mu_1 \geq \mu_2 \geq \cdots > 0$，
对应的标准正交特征向量 $(u_n)$；$R(x) = \frac{\langle x,
Ax\rangle}{\norm x^2}$。极小–极大公式见
\cref{exo:b3:spectral:8}；这里加以应用。
\begin{enumerate}[resume]
  \item （幂法）对 $x \neq 0$，记 $m_p =
        \sum_n\mu_n^p\abs{\langle u_n, x\rangle}^2$。证明
        $m_pm_{p+2} \geq m_{p+1}^2$（由柯西–施瓦茨不等式），由此推出
        如下不等式链：
        \[
        R(x) \;\leq\; \frac{\langle Ax, Ax\rangle}{\langle x,
        Ax\rangle} \;\leq\; R(Ax) \;\leq\; \mu_1,
        \]
        并证明：若 $\langle u_1, x\rangle \neq 0$，
        则 $R(A^kx) \to \mu_1$。迭代地将算子作用于
        任意一般向量，即可算出最大特征值 ---
        这正是数值分析中的幂法，并附有证明。
  \item （外尔稳定性）对紧自伴正算子
        $A, B$，由\cref{exo:b3:spectral:8}推出
        \[
        \abs{\mu_n(A) - \mu_n(B)} \;\leq\; \vertiii{A - B}
        \qquad\text{对每个 } n :
        \]
        整个谱在算子
        范数下是 $1$-利普希茨的 ---
        大型对称系统的特征值，可由带有误差保证的近似值算出。
  \item 对 $G$ 应用瑞利界，取检验
        函数 $u(x) = x(1-x)$，解 $-w'' = u$，$w(0) =
        w(1) = 0$，得 $Gu = w = \frac{x(1-x)(1 + x -
        x^2)}{12}$，计算
        \[
        \norm u_2^2 = \frac1{30}, \qquad
        \langle u, Gu\rangle = \frac1{12}\Bigl(\frac1{30} +
        \frac1{140}\Bigr) = \frac{17}{5040},
        \qquad R(u) = \frac{17}{168},
        \]
        由此得到\emph{严格}的界 $\frac1{\pi^2}
        = \lambda_1 \geq \frac{17}{168}$，即 $\pi \leq
        \sqrt{168/17} < 3.1437$。
  \item 取第 11 题不等式链中的一步，对同一检验
        函数：利用 $\int_0^1(x - x^2)^4\dd x =
        \frac1{630}$，计算
        \[
        \norm{Gu}_2^2 = \frac1{144}\Bigl(\frac1{30} +
        \frac2{140} + \frac1{630}\Bigr) = \frac{31}{90720},
        \qquad
        \frac{\langle Gu, Gu\rangle}{\langle u, Gu\rangle} =
        \frac{31}{306},
        \]
        由此得出 $\pi \leq \sqrt{306/31} < 3.1419$：两次
        积分，得到四位正确数字。（每多迭代一次，
        精度大致平方一次：
        特征向量的间隙 $\lambda_1/\lambda_2 = 4$
        导致几何收敛。）
\end{enumerate}

\medskip
\noindent\textbf{第六部分 --- $G^2$ 的迹与
$\zeta(4)$。}
\begin{enumerate}[resume]
  \item 证明 $\sum_n\lambda_n^2 =
        \iint_{\intcc01^2}g(x,y)^2\,\dd x\,\dd y$
        \emph{（将 $g$ 按乘积基
        $(e_m(x)e_n(y))_{m,n}$ 展开，该基是 $L^2(\intcc01^2)$ 中的 ---
        希尔伯特基，参见\cref{exo:b3:spectral:5} --- 并
        在正方形上应用帕塞瓦尔恒等式；第 7 题确定了
        这些系数）}。
  \item 计算二重积分：
        \[
        \iint g^2 = 2\int_0^1(1-x)^2\Bigl(\int_0^x
        y^2\,\dd y\Bigr)\dd x
        = \frac23\int_0^1x^3(1-x)^2\,\dd x = \frac1{90} .
        \]
  \item 得出 $\zeta(4) = \dfrac{\pi^4}{90}$，并
        不作计算，说明高次幂
        $G^k$ 的迹如何产生 $\zeta(2k) \in \pi^{2k}\,\Q$，对所有
        $k \geq 1$ 成立，并说明为何奇数值 $\zeta(3),
        \zeta(5), \dots$ 在结构上超出这台
        机器的能力范围。
  \item （$\pi$ 的下界）由 $\lambda_1^2 \leq
        \sum_n\lambda_n^2 = \frac1{90}$，可得 $\pi \geq
        90^{1/4} > 3.080$，并结合第 14 题，给出
        双侧结论
        \[
        3.080 \;<\; \pi \;<\; 3.1419,
        \]
        这一结论完全来自振动弦的
        算术。若改用更高次的迹
        $(\operatorname{tr}G^{2k})^{-1/4k}$，哪一侧收敛更快，
        以及为什么？
\end{enumerate}

\medskip
\noindent\textbf{第七部分 --- 受迫与共振。}设
$\nu \in \R$，考虑受迫弦 $-u'' - \nu u =
f$，$u(0) = u(1) = 0$，$f \in L^2$ 且 $c_n = \langle
e_n, f\rangle$。
\begin{enumerate}[resume]
  \item 设 $\nu \notin \{n^2\pi^2 : n \geq 1\}$，证明
        下式
        \[
        u = \sum_{n\geq1}\frac{c_n}{n^2\pi^2 - \nu}\,e_n
        \]
        在 $L^2$ 中收敛，且在 $\intcc01$ 上一致收敛
        \emph{（由 $(c_n)$ 与
        $\sum n^{-4}$ 的尾部之间的柯西–施瓦茨不等式，并取 $\norm{e_n}_\infty =
        \sqrt2$）}，并证明它满足 $u = Gf + \nu Gu$
        --- 弗雷德霍姆择一性的显式坐标形式，
        （见\cref{thm:b3:spectral:fredholm}），附带
        唯一性。
  \item 设 $\nu = m^2\pi^2$。证明方程 $u = Gf + \nu
        Gu$ 有解当且仅当 $c_m = 0$，解唯一，至多相差
        $e_m$ 的倍数。物理意义：
        恰以秋千自身的固有频率推动它。
  \item 设 $\nu < \pi^2$，证明解算子
        $R_\nu\colon f \mapsto u$ 在 $L^2$ 上有界，
        范数为 $\frac1{\pi^2 - \nu}$，紧、自伴，
        且为正算子：整个谱分析重新开始，
        只是谱整体平移了 $\nu$。
  \item （综合）整理本题的对应关系：
        特征值 $\leftrightarrow$ 频率的平方
        （谐波）；迹 $\leftrightarrow$ $\zeta(2)$；
        希尔伯特–施密特范数 $\leftrightarrow$ $\zeta(4)$；
        弗雷德霍姆择一性 $\leftrightarrow$ 共振；
        极小–极大 $\leftrightarrow$ 变分界（$\pi <
        3.1437$ 来自一个多项式）。一个积分算子，
        由此兑现了五章分析的成果。
\end{enumerate}

\medskip
\noindent\textbf{第八部分 --- 最后三则回响。}
\begin{enumerate}[resume]
  \item （$\zeta(6)$，白得）第 6 题的帕塞瓦尔
        恒等式给出 $\sum_{n\text{ 奇数}}n^{-6} =
        \frac{\pi^6}{960}$。将 $\zeta(6)$ 拆分为奇数与
        偶数 $n$ 两部分，并得出
        \[
        \zeta(6) = \frac{\pi^6}{945},
        \]
        无需新的积分：第 17 题的机器，
        （取 $G^3$ 的迹）本可得到同一数值，
        但需付出迭代核的代价 --- 在一个选得恰当的函数上
        应用帕塞瓦尔恒等式，是此处更省事的途径。
  \item （基态为正）设 $A$ 为
        $L^2(\intcc01)$ 上的紧自伴正算子，
        由连续对称核 $k > 0$ 给出，它定义在
        $\intoo01^2$ 上，最大特征值为 $\mu_1$。证明：
        （a）瑞利商的任一极大点都是
        $\mu_1$ 对应的特征函数；（b）若 $u$ 为其一，则
        $\langle\abs u, A\abs u\rangle \geq \langle u,
        Au\rangle$，且当 $u$ 在正测度集上取两种符号时，为\emph{严格}不等式 ---
        因此 $u$
        几乎处处保持定号，且 $u = \mu_1^{-1}Au$
        在 $\intoo01$ 上处处不为零；（c）$\mu_1$ 是
        \emph{单}特征值。在 $G$ 上逐条验证：
        $e_1 = \sqrt2\sin(\pi x) > 0$，而对每个 $e_n$，$n
        \geq 2$，由于与 $e_1$ 正交，必须变号
        （事实如此：有 $n - 1$ 个内部零点）。
  \item （到谱的距离，以及共振的
        代价）。对 $\nu \notin \{n^2\pi^2\}$，证明
        第 19 题中的解算子 $R_\nu$ 是
        有界、自伴、紧，并且满足
        \[
        \vertiii{R_\nu} =
        \frac1{\min_{n\geq1}\,\abs{n^2\pi^2 - \nu}}
        = \frac1{\operatorname{dist}\bigl(\nu,
        \{n^2\pi^2\}\bigr)},
        \]
        范数在最近的模式上取到。于是
        量化第 20 题中秋千的情形：以 $f = e_1$ 作为强迫项，
        在 $\nu = (1 - \varepsilon)\pi^2$ 处得到 $u =
        \frac{1}{\varepsilon\pi^2}\,e_1$，相当于
        $\frac1\varepsilon$ 倍地放大静态响应 $Ge_1 =
        \frac1{\pi^2}e_1$ --- 比基频低百分之一
        （$\varepsilon = 10^{-2}$）时，弦的
        响应增强了一百倍。
\end{enumerate}
\end{problem}
